% ---------------------------------------------------------------------------
% Stress-aware grasp synthesis.
%
% Replaces \subsection{Stress-aware grasp synthesis} (sec:method:synth) in method.tex.
% Written against the pipeline that actually generated the data: smgrasp/finger_grasp_final.py
% driven by grasp_synthesis/collect_demos_synth_v4.py. The earlier draft described
% smgrasp/finger_grasp.py (the v3 collector), which differs in the objective, the selection
% rule and the derived-parameter claims.
%
% The appendix block at the end holds the numeric detail; move it behind \appendix.
% Algorithm 1 is built from base LaTeX (no algorithm/algpseudocode in this install).
% ---------------------------------------------------------------------------

\subsection{Stress-aware grasp synthesis}
\label{sec:method:synth}
For each object instance the demonstrator plans a $7$-DoF grasp
$x=(\mathbf{t},\mathbf{r},\omega)$: tool-centre position, orientation, and \emph{commanded}
gripper width. Planning runs on the mesh, outside the physics simulator.

The workflow is three parts. A width-controlled finite-element contact model predicts, for a
candidate grasp, the stress it would impose and the force it would hold with; the model is
factored once per geometry, which makes both the search and randomisation over material
properties cheap. A feasibility ladder then excludes candidates that miss the object, penetrate it,
or cannot carry it, and among those that survive an objective trades bulk deformation against
local contact pressure. Finally a search --- geometric seeding, a batched
scoring pass, local optimisation, and a width rescan --- returns the gentlest holdable grasp,
re-run under progressively looser bounds if the object admits none. We detail each in turn.

\noindent\textbf{Width-controlled contact model.}
Common parallel jaw grippers are position-controlled: one commands a width and the grip force is an
outcome. We model this directly. The object is tetrahedralised and a linear-elastic stiffness
matrix assembled; since it is free-floating the system is singular in the six rigid-body modes,
so we solve a bordered inertia-relief system, factored once per geometry. Two rigid pads, whose
rectangular footprints are taken from the real gripper's finger meshes, are closed to width
$\omega$. Boundary nodes inside a footprint are pushed onto the pad's flat plane and constrained
\emph{only} along the closing axis, leaving tangential motion free so the material bulges; a
thin smoothstep fillet at the pad edge avoids a flat-punch stress singularity. Contact is
normal-only, and friction enters solely through the holdability test below.

\noindent\textbf{Material scaling.}
Because the displacement is prescribed rather than the force, the deformation field is
independent of Young's modulus, instead is dictated by the grasp width. Thus below yield both stress and grip force are linear in it,
\begin{equation}
  \sigma(E) = E\,\sigma_1, \qquad F(E) = E\,F_1 .
  \label{eq:E-linear}
\end{equation}
One solve at $E\!=\!1$ therefore yields every material by scalar multiplication, which makes
randomising stiffness across the demonstration set essentially free. A grasp is
\emph{holdable} when Coulomb friction at the two pads carries the object's weight with an
acceleration margin, $2\mu F(E) \ge m\,(g+a)$; we use $\mu\!=\!0.7$ and $a\!=\!g$, a quasi-static
$2g$ criterion rather than a wrench-closure analysis.


\begin{table}[t]\centering\footnotesize
\caption{Feasibility gates. $\delta$ is the pad indentation, $F$ the per-jaw normal force,
$A_{\min}$ the smaller pad's contact area, $\hat{\mathbf{a}}$ the closing axis,
$\boldsymbol{\tau}$ the gravitational torque about the object's centre of mass, and $\sigma_Y$ the
material's yield stress. A candidate failing a gate returns a penalty shaped by its margin of
failure, so the search retains a gradient toward the feasible band rather than a flat infeasible
plateau. Nominal bounds are given in Appendix~\ref{app:synth:constants}.}
\label{tab:gates}
\begin{tabular}{@{}ll@{}}
\toprule
gate & feasible when \\
\midrule
table       & fingers stay above the table plane \\
penetration & finger bodies stay outside the object \\
contact     & each pad captures a non-empty contact patch \\
indentation & $\delta \le \delta_{\max}$ (linear-elastic regime) \\
holdability & $2\mu F \ge m\,(g+a)$ \\
torsion     & $\bigl|\boldsymbol{\tau}\cdot\hat{\mathbf{a}}\bigr| \le
               \tfrac{4}{3}\,\mu F \sqrt{A_{\min}/\pi}$ \\
yield       & $\sigma_{\mathrm{blk}} \le \sigma_Y$ \\
\bottomrule
\end{tabular}
\end{table}

\noindent\textbf{Objective.}
A candidate is \emph{feasible} when it clears every gate of Table~\ref{tab:gates}: the fingers
clear the table and the object's interior, both pads make contact at an indentation the
linear-elastic model still supports, and the grasp holds the object against sliding, twisting and
yielding. Among feasible candidates we maximise
\begin{equation}
  S(x) \;=\; -\,\sigma_{\mathrm{blk}} \;-\; w_{\mathrm{pr}}\,\frac{F}{A_{\min}},
  \label{eq:score}
\end{equation}
where $\sigma_{\mathrm{blk}}$ be the mean of the top decile of element stresses with contact-adjacent elements masked out, so that it reports
\emph{bulk} deformation rather than the mesh-dependent contact singularity.  $A_{\min}$ is the contact area of the worse pad, so the second term is the mean contact
pressure. 
% The pressure term is what separates enveloping grasps from stem-type ones: squeezing a
% stiff, slender part deforms the body \emph{less} and therefore scores better on bulk stress alone,
% but concentrates the same holding force on a far smaller area. 
Bulk gentleness and local gentleness are thus both represented, and neither is a hand-set floor.

\noindent\textbf{Search.}
The search is a funnel (Algorithm~\ref{alg:synth}). Candidates are seeded from two geometric
generators --- antipodal surface pairs inside the friction cone, and medial-axis points closing
across the local tangent. Each primitive determines a complete $7$-DoF pose: its axis and anchor
fix the orientation and position, and the object's local cross-section there fixes the commanded
width (Appendix~\ref{app:synth:seeds}).

Seeds are pre-screened on the two gates that batch cheaply --- table clearance and finger
penetration, at a deliberately loose tolerance --- and on an orientation box. Survivors are
scored in one batched solve, where the full ladder of Table~\ref{tab:gates} applies, and the
best few initialise a multi-start CMA-ES run whose proposals are scored the same way. The
orientation box bounds the search rather than gating it: roll and pitch are kept away from
extreme tilts, at which the wrist and forearm would strike the table --- geometry the
finger-level gates do not see --- and yaw is limited so the fingers do not occlude the object in
the very view the policy must act on. Poses outside it are never evaluated, whereas a candidate
failing a gate is evaluated and penalised, and it is the bound the relaxation ladder below
widens first. The CMA runs search all seven degrees of freedom jointly and advance in lockstep, so each
generation is a single batched solve. A final pass then holds each of the best spatially
distinct poses \emph{fixed} and sweeps the one remaining degree of freedom, width, exhaustively
on a grid several times finer than the optimiser's step, keeping the widest holdable
value. Width is singled out because it is what the real gripper actually commands, and because
the joint search resolves it least well.

\noindent\textbf{Relaxation ladder.}
A single fixed feasibility band cannot serve an object set spanning three orders of magnitude in
stiffness and scale: bounds loose enough for the hardest object admit poor grasps on the
easiest. The funnel is therefore run inside a short ladder. The first pass uses nominal bounds;
if it yields no holdable grasp, a second pass widens the orientation box and the penetration
tolerance. The candidates this admits skew toward pinch grasps, so the same pass \emph{doubles}
the pressure weight of Eq.~\eqref{eq:score}: the constraints loosen and the objective tightens
together, buying reach without buying grasps that hold an object on two small patches. A third
pass additionally releases the yield gate, accepting a grasp that marks the object rather than
failing to lift it. This last release is a deliberate trade: a demonstration that never lifts
teaches the policy nothing, whereas one that marks the object still carries the approach, grasp
and lift structure we need from it. It is a last resort rather than a licence: $2.9\%$ of
planned grasps released the gate (Appendix~\ref{app:synth:tiers}). If none succeeds the planner
returns its most-centred feasible seed and the episode is recorded as an attempt. The ladder
stops at the first pass that produces a holdable grasp, so it costs nothing on objects that do
not need it.

\begin{figure}[t]
\hrule\vspace{3pt}
\noindent\textbf{Algorithm~1}~~Stress-aware grasp synthesis
\vspace{2pt}\hrule\vspace{4pt}
\footnotesize
\noindent\begin{tabular}{@{}r@{~~}p{0.86\linewidth}@{}}
   & \textbf{Input:} mesh, pose, material $(E,\rho,\mu)$, table height \\
   & \textbf{Output:} grasp $x=(\mathbf{t},\mathbf{r},\omega)$ and its tier \\[2pt]
 1 & assemble FEM, factor the bordered system \emph{once} \\
 2 & $P \leftarrow$ antipodal pairs $\cup$ medial-axis points \\
 3 & $X_0 \leftarrow$ \textsc{LiftToPose}($P$) \hfill\emph{axis, height, local cross-section} \\[3pt]
 4 & \textbf{for} tier $t = 0,1,2$ \textbf{do} \\
 5 & \quad bounds $\leftarrow$ \textsc{Relax}($t$) \\
 6 & \quad $X_1 \leftarrow \{x \in X_0 :$ table, orientation, penetration $\}$ \\
 7 & \quad score $X_1$ in one batched solve; $X_2 \leftarrow$ best $K$ \\
 8 & \quad $X_3 \leftarrow$ \textsc{CMA-ES}($X_2$) over all $7$ DoF, runs in lockstep \\
 9 & \quad $X_4 \leftarrow$ spatially distinct poses of $X_3$ \\
10 & \quad $x^\ast \leftarrow \arg\max S(x)$ over a $1$-DoF width sweep at each pose of $X_4$ \\
11 & \quad \textbf{if} $x^\ast$ is holdable \textbf{then return} $(x^\ast, t)$ \\[3pt]
12 & \textbf{return} most-centred feasible seed, tier $3$ \\
\end{tabular}
\vspace{4pt}\hrule
\caption{The funnel (lines 6--10) is re-run under progressively looser bounds and returns at the
first tier that yields a holdable grasp.}
\label{alg:synth}
\end{figure}

% ===========================================================================
%  APPENDIX MATERIAL --- move behind \appendix
% ===========================================================================

\section{Grasp synthesis: implementation detail}
\label{app:synth}

\subsection{Seed construction}
\label{app:synth:seeds}
The closing axis fixes the tool $y$ axis, with the approach taken as near vertical as that axis
permits and yaw folded into $[-90^\circ,90^\circ]$ by jaw symmetry. The pads are centred on the
anchor, the fingertip height is sampled between table clearance and the object's top, and the
commanded width is the object's span inside the pad footprint at that height, less a small closure
margin. The span is measured as below.

\subsection{Local cross-section on objects with holes}
\label{app:synth:width}
A seed's commanded width is the object's thickness at the anchor along the intended closing axis.
The natural estimate collects surface vertices in a thin cylinder around the anchor aligned with
that axis, projects them onto it, and takes the extent. On a convex body this is exactly the local
thickness.

On a ring, a handle or a deep concavity it is not: the cylinder passes through the void and
collects surface on the far side, so the extent measures the whole object rather than the material
between the fingers. On our smaller torus --- $100$~mm outer diameter, $19$~mm tube --- the
estimate returns $100$~mm where the correct value is $19$~mm. Every medial seed therefore asked the
gripper to close on the entire ring; $63\%$ saturated at the $79$~mm jaw limit, and none of those
attempts lifted the object.

We instead march the signed distance field outward from the anchor along $\pm$ the closing axis in
$0.5$~mm steps and take the first exit in each direction; their sum is the \emph{solid run}
actually traversed. The solid run replaces the extent only when it is below a quarter of it, that
is, only when a genuine void was crossed. On a convex body the two agree and the rule is inert, so
it changes rings and handles and nothing else. Demonstrator success on the small torus rose from
$13.8\%$ to $47.2\%$.

\subsection{Statistics computed but not scored}
\label{app:synth:gates}
An alignment statistic (the mean absolute cosine between the closing axis and the local surface
normals) and an unmasked $98$th-percentile stress are computed and logged for analysis but do
\emph{not} enter Eq.~\eqref{eq:score}.

\subsection{Relaxation tiers}
\label{app:synth:tiers}
\begin{center}\footnotesize
\begin{tabular}{@{}cp{0.62\linewidth}@{}}
\toprule
tier & settings \\
\midrule
$0$ & nominal: roll $\pm25^\circ$, pitch $\pm20^\circ$, yaw $\pm60^\circ$ ($\pm90^\circ$ for objects
      whose largest extent exceeds $70$~mm, which cannot be hidden by the fingers at any yaw);
      penetration $5$~mm; indent $10$~mm \\
$1$ & roll/pitch $+10^\circ$, yaw $\pm80^\circ$, penetration $20$~mm, indent $12$~mm; yield gate
      kept. Pressure weight $\times 2$ --- a tightening, not a relaxation \\
$2$ & as tier $1$, yield gate released \\
$3$ & no tier succeeded: the feasible seed nearest the object centre of mass \\
\bottomrule
\end{tabular}
\end{center}

The planner records the tier that produced each grasp. Over $12{,}617$ planning attempts spanning
$291$ of the $294$ collection runs:

\begin{center}\footnotesize
\begin{tabular}{@{}lrrrr@{}}
\toprule
object group & tier $0$ & tier $1$ & tier $2$ & fallback \\
\midrule
all                 & $93.2\%$ & $3.0\%$  & $2.9\%$  & $0.9\%$ \\
original subset     & $93.1\%$ & $2.5\%$  & $3.9\%$  & $0.5\%$ \\
primitive campaign  & $99.8\%$ & $0.2\%$  & $0.0\%$  & $0.0\%$ \\
rod and tori        & $62.9\%$ & $21.4\%$ & $12.1\%$ & $3.6\%$ \\
held-out categories & $77.2$--$95.6\%$ & $3.3$--$7.2\%$ & $0.7$--$7.8\%$ & $0.4$--$7.7\%$ \\
\bottomrule
\end{tabular}
\end{center}

The spread is the point. The primitive-shape campaign essentially never needs the ladder, while the
rod-and-torus group --- chosen because a thin rod and a ring are hard to grasp gently --- needs it
on a third of attempts. Rates are per attempt; weighting each object or each category equally
instead moves the tier-$2$ share by about a point. Attempts exceed episodes because a plan is made
for every attempt, including those that go on to fail.

\subsection{Constants}
\label{app:synth:constants}
Gate bounds (Table~\ref{tab:gates}): $2$~mm table tolerance, $5$~mm finger-body penetration,
$\delta_{\max}=10$~mm indentation, $\mu=0.7$ and $a=g$ in the holdability test, and the object's
own $\sigma_Y$. Seeds: $2600$ antipodal pairs and $500$ medial-axis points, the latter expanded to
four widths each. Filter: $10$~mm per-finger penetration. Search: $K=6$ seeds carried forward, $400$
evaluations per seed, poses counted distinct beyond $5$~mm or $10^\circ$, the best $15$ carried to
the width rescan, which samples $25$ widths over $\pm3$~mm at $0.25$~mm resolution. Gripper range
$8$--$79$~mm; Poisson ratio $0.33$; $\sim\!1500$ tetrahedra per object. These are fixed across the
entire object set; none is set per object or per category. The one object-dependent bound is the
yaw limit, which switches at a $70$~mm extent as described above.
